% Draft (2026-11-07) for the ICML 2027 version. \tbd{} marks numbers from the preregistered held-out runs.
\section{Introduction}
\label{sec:intro}

MBPP+ problem 102 \citep{liu2023evalplus} asks for a function that converts a snake-case string to camel case and
gives three ordinary examples. A generated function splits on underscores, capitalises each word and passes
all three. The hidden suite then passes \verb|___python_program| and expects \verb|___PythonProgram|; the function
returns \verb|PythonProgram|. The task says nothing about how to handle leading underscores. The benchmark treats
this program as defective. By that measure, an automated auditor that reads and passes it has missed a defect.

Models increasingly check what other models write. Audit benchmarks score them as in this example, against a hidden
reference. We argue that these scores conflate two kinds of case. In some, the evidence available to the auditor
determines the correct behaviour, and a miss is the auditor's failure. In others, such as the example above, the
evidence is equally consistent with the program and the reference; where the evidence is moreover identically
distributed under both behaviours, as our construction guarantees, no auditor, whatever its capacity, can tell them
apart without being told which behaviour is wanted. A ``detection'' in these cases amounts to agreement with the
benchmark's convention; it would be a false accusation had the convention gone the other way.

This distinction is classical. For any flagging rule, recall minus false-flag rate is bounded by the total
variation between the evidence a defective and a correct program produce. The best attainable recall at a given
false-flag rate is the Neyman--Pearson power of that evidence (Section~\ref{sec:framework}). Audit evaluations have
lacked a way to observe this distinction because a benchmark fixes one contract and hides the counterfactual.
We construct that counterfactual experimentally (Section~\ref{sec:e6}). Each unit starts with a real benchmark task
that leaves one behaviour open; we build two programs that differ only in that behaviour. We grade the auditor against a
contract drawn independently of what it sees, then reveal the contract in different ways. Execution validates the
programs, oracles and boundaries of the open behaviour; two models, which also implement the rules, judge whether both are admissible.

\textbf{Findings.} Across \NSixUnits{} held-out units and auditors from two vendors at two reasoning budgets:
\begin{itemize}\itemsep1pt
\item With the contract withheld, flagging cannot track which program is wrong: the same readings are scored under
  both contracts, so recall equals the false-flag rate (\NSixNoneFlag\%) by construction. The informative difference
  is which program is flagged. The one that departs from the benchmark's behaviour is flagged \NSixFlagPref{} points
  [\NSixFlagPrefLo, \NSixFlagPrefHi] more often than the one that follows it, and diagnosed more often too (the registered convention advantage, half
  the difference between the two diagnosis rates, is \NSixTwo{} points [\NSixTwoLo, \NSixTwoHi]).
\item The models' own solutions written from the task alone follow the benchmark's behaviour on these units 89\%
  of the time. The convention the benchmark rewards is one the models share, and rewarding it would be a false
  accusation had the task's author meant the other rule.
\item Adding the missing requirement as one ordinary sentence raises verified diagnosis by \NSixFive{} points
  [\NSixFiveLo, \NSixFiveHi] over the task alone, with no correct program flagged in that arm (\NSixResolveFpr\% when the contract is
  stated as a choice between the two rules).
\item When told only that the task leaves the choice open, the two vendors respond differently: one passes almost
  every program, while the other asks for clarification.
\item On positive controls, where the task's words rule out one program, the same auditors separate the programs
  without being told (recall minus false-flag rate \NSixFourJ{} points).
\end{itemize}
Among natural hidden-test failures (Section~\ref{sec:natural}), fewer than a quarter are ones that both classifiers
agree the task's words determine. Three earlier preregistered studies of scientific code, results and data
(Section~\ref{sec:transfer}) are consistent with the same distinction, although none of them varied the information
an auditor receives in the way our experiment does.

\textbf{Contributions.} (i) A randomised-contract design that makes the information limit of an audit observable and
separates convention agreement, unjustified accusation, requests for clarification and verified diagnosis, with
execution-validated programs and oracles. (ii) Preregistered evidence that, on such cases, auditors side with the benchmark's convention, and that one
added sentence stating the requirement turns that agreement into verified diagnosis. (iii) A description of
natural hidden-test failures by whether the task's words determine the expected value, and an
interpretation of scientific code, results and data audits using the same distinction between capacity and information misses.
Code, records and preregistrations are released.
